\subsection{System Overview}

The proposed system implements a bidirectional haptic interaction loop consisting of four principal layers: the haptic device hardware, embedded control firmware, host-side simulation and haptic rendering, and graphical visualisation. A conceptual overview of the signal and force flow is shown in Fig.~\ref{fig:system_architecture}.

\begin{figure}[!ht]
	\centering
	\resizebox{1\textwidth}{!}{%
		\begin{circuitikz}
			\tikzstyle{every node}=[font=\fontsize{11.9pt}{15.5pt}\selectfont]
			\draw [ rounded corners = 6.0] (2.125,13.375) rectangle (4.25,12.875);
			\node[anchor=center, align=center, inner xsep=0.080cm, inner ysep=0.085cm, rounded corners=0.020cm] at (3.1875,13.125) {\fontsize{11.9pt}{15.5pt}\selectfont User};
			\draw [ rounded corners = 6.0] (11.25,13.5) rectangle (14.25,12.75);
			\node[anchor=center, align=center, inner xsep=0.080cm, inner ysep=0.085cm, rounded corners=0.020cm] at (12.75,13.125) {\fontsize{11.9pt}{15.5pt}\selectfont Host Interface};
			\draw [ rounded corners = 6.0] (15,11.375) rectangle (18,10.625);
			\node[anchor=center, align=center, inner xsep=0.080cm, inner ysep=0.085cm, rounded corners=0.020cm] at (16.5,11) {\fontsize{11.9pt}{15.5pt}\selectfont visualisation};
			\draw [ rounded corners = 6.0] (14.75,15.625) rectangle (18.375,14.875);
			\node[anchor=center, align=center, inner xsep=0.080cm, inner ysep=0.085cm, rounded corners=0.020cm] at (16.5625,15.25) {\fontsize{11.9pt}{15.5pt}\selectfont Haptic Rendering};
			\draw [ rounded corners = 6.0] (18.625,13.5) rectangle (22.625,12.75);
			\node[anchor=center, align=center, inner xsep=0.080cm, inner ysep=0.085cm, rounded corners=0.020cm] at (20.625,13.125) {\fontsize{11.9pt}{15.5pt}\selectfont Physics Simulation};
			\draw [-{Stealth[scale=1.5]}, ] (4.125,13.375) arc[start angle=121.85, end angle=64.88, radius=2.2317]node[pos=0.48,above, fill={rgb,255:red,255; green,255; blue,255}, fill opacity=1, text opacity=1, inner xsep=0.080cm, inner ysep=0.085cm, rounded corners=0.020cm]{motion};
			\draw [ rounded corners = 6.0] (6.125,13.5) rectangle (9.125,12.75);
			\node[anchor=center, align=center, inner xsep=0.080cm, inner ysep=0.085cm, rounded corners=0.020cm] at (7.625,13.125) {\fontsize{11.9pt}{15.5pt}\selectfont Haptic Device};
			\draw [-{Stealth[scale=1.5]}, ] (6.25,12.75) arc[start angle=-59.83, end angle=-126.90, radius=1.9265]node[pos=0.48,below, fill={rgb,255:red,255; green,255; blue,255}, fill opacity=1, text opacity=1, inner xsep=0.080cm, inner ysep=0.085cm, rounded corners=0.020cm]{haptic feedback};
			\draw [-{Stealth[scale=1.5]}, ] (9,13.5) arc[start angle=122.06, end angle=57.94, radius=2.2368]node[pos=0.48,above, fill={rgb,255:red,255; green,255; blue,255}, fill opacity=1, text opacity=1, inner xsep=0.080cm, inner ysep=0.085cm, rounded corners=0.020cm]{state packets};
			\draw [-{Stealth[scale=1.5]}, ] (11.375,12.75) arc[start angle=-52.24, end angle=-127.76, radius=1.9390]node[pos=0.48,below, fill={rgb,255:red,255; green,255; blue,255}, fill opacity=1, text opacity=1, inner xsep=0.080cm, inner ysep=0.085cm, rounded corners=0.020cm]{torque commands};
			\draw [-{Stealth[scale=1.5]}, ] (12.875,13.5) arc[start angle=-173.12, end angle=-280.83, radius=1.5881]node[pos=0.3,above, fill={rgb,255:red,255; green,255; blue,255}, fill opacity=1, text opacity=1, inner xsep=0.080cm, inner ysep=0.085cm, rounded corners=0.020cm]{joint state};
			\draw [-{Stealth[scale=1.5]}, ] (15.25,14.875) arc[start angle=24.32, end angle=-83.81, radius=1.2447]node[pos=0.15,below, fill={rgb,255:red,255; green,255; blue,255}, fill opacity=1, text opacity=1, inner xsep=0.080cm, inner ysep=0.085cm, rounded corners=0.020cm]{force command};
			\draw [{Stealth[scale=1.5]}-, ] (15,11) -- (13,12.75)node[pos=0.31,above, fill={rgb,255:red,255; green,255; blue,255}, fill opacity=1, text opacity=1, inner xsep=0.080cm, inner ysep=0.085cm, rounded corners=0.020cm]{joint state};
			\draw [-{Stealth[scale=1.5]}, ] (17.125,14.875) -- (17.125,11.375)node[pos=0.73,above, fill={rgb,255:red,255; green,255; blue,255}, fill opacity=1, text opacity=1, inner xsep=0.080cm, inner ysep=0.085cm, rounded corners=0.020cm]{proxy state};
			\draw [-{Stealth[scale=1.5]}, ] (18.25,14.875) -- (19.125,13.5)node[pos=0.85,above, fill={rgb,255:red,255; green,255; blue,255}, fill opacity=1, text opacity=1, inner xsep=0.080cm, inner ysep=0.085cm, rounded corners=0.020cm]{interaction force};
		\end{circuitikz}
	}%
	\caption{High-level architecture of the proposed haptic system. User motion is measured by the haptic device and transmitted to the host interface. The host distributes state information to the haptic rendering and visualisation subsystems, while the haptic rendering module computes interaction forces, updates the simulation, and generates feedback commands that are returned to the device as actuator torque commands.}
	\label{fig:system_architecture}
\end{figure}

User motion is measured at the device joints and transmitted to the host in real time. The host computes interaction forces based on the virtual environment state, maps these to joint torque commands, and returns them to the device actuators, forming a closed-loop interaction.

To support stable real-time operation, time-critical control tasks are executed on the embedded system at high frequency, while simulation, rendering, and visualisation are performed asynchronously on the host.

\subsection{Design Constraints and Assumptions}

The system is designed to prioritise real-time stability, modularity, and reproducibility within the constraints of a low-cost prototype.

Stable haptic interaction requires high update rates (approximately 1\,kHz) and low end-to-end latency. Accordingly, the architecture assumes high-frequency embedded torque control and asynchronous host-side computation.

Mechanical and control limits are selected to ensure safe operation while maintaining sufficient force capability for meaningful interaction.

\subsection{Hardware Architecture}

The haptic interface is implemented as a planar two-degree-of-freedom mechanism with orthogonal revolute joints. Each joint is actuated by a brushless motor operating in torque control mode and instrumented with absolute angular position sensing.

The structure is designed to be lightweight and back-drivable to minimise reflected inertia and improve transparency.

\subsection{Embedded Control Architecture}

The embedded system performs real-time actuator control, sensor acquisition, and communication with the host. Joint states are streamed to the host at a fixed rate, while torque commands are received and applied continuously.

Safety mechanisms, including torque limiting and communication watchdogs, are implemented to ensure predictable behaviour under fault conditions.

\subsection{Communication Architecture}

Communication between the embedded controller and host is implemented using a packet-based interface exchanging joint state and torque command data.

The protocol is designed to support deterministic timing and low overhead suitable for real-time haptic interaction.

\subsection{Host Software Architecture}

The host system is implemented as a modular multi-threaded application comprising subsystems for world management, physics simulation, haptic rendering, device communication, and visualisation.

These subsystems operate at different rates and exchange data through structured communication channels, allowing high-frequency haptic computation to be maintained independently of slower processes such as rendering.

\subsection{Haptic Rendering Paradigm}

Interaction forces are generated using a proxy-based rendering approach in which the device position is coupled to a constraint-based proxy within the virtual environment. This enables stable rendering of contact interactions under discrete-time computation.

\subsection{Force Transmission to the Device}

Cartesian interaction forces are mapped to joint torques using the Jacobian transpose formulation:

\begin{equation}
	\tau = J^{T}(q)F
\end{equation}

where $\tau$ is the joint torque vector, $J(q)$ is the manipulator Jacobian, and $F$ is the rendered force.

\subsection{Visualisation Layer}

A graphical visualisation subsystem provides real-time feedback of the simulated environment. Rendering is decoupled from the haptic loop to prevent graphical latency from affecting interaction stability.